\documentclass[10pt,letterpaper]{article}
\usepackage[utf8]{inputenc}
\usepackage[T1]{fontenc}
\usepackage[margin=0.35in]{geometry}
\usepackage{xcolor}
\usepackage[hidelinks]{hyperref}

\definecolor{primary}{RGB}{18, 38, 68}
\definecolor{textdark}{RGB}{30, 30, 30}
\definecolor{linkcolor}{RGB}{0, 75, 145}

\hypersetup{colorlinks=true, urlcolor=linkcolor, linkcolor=linkcolor}
\pagestyle{empty}
\raggedbottom
\raggedright

\newcommand{\resumesection}[1]{%
  \vspace{2pt}%
  {\color{primary}\normalsize\bfseries\MakeUppercase{#1}}%
  \vspace{1pt}\hrule height 0.5pt\vspace{2pt}%
}

\def\labelitemi{\color{primary}$\bullet$}
\renewenvironment{itemize}{%
  \begin{list}{\color{primary}$\bullet$}{%
    \setlength{\leftmargin}{1.1em}%
    \setlength{\itemsep}{0.2pt}%
    \setlength{\topsep}{0.5pt}%
    \setlength{\parsep}{0pt}%
    \setlength{\partopsep}{0pt}%
  }%
}{%
  \end{list}%
}

\newenvironment{skilllist}{%
  \begin{list}{}{%
    \setlength{\leftmargin}{0pt}%
    \setlength{\itemsep}{0.2pt}%
    \setlength{\topsep}{0.5pt}%
    \setlength{\parsep}{0pt}%
    \setlength{\partopsep}{0pt}%
  }%
}{%
  \end{list}%
}

\begin{document}
\color{textdark}

\begin{center}
    {\LARGE \textbf{\color{primary} Alex Rivera}} \\ \vspace{1.5pt}
    {\large \textbf{Analista de Datos \textbar\ Machine Learning Aplicado, SQL Avanzado \& Modelado Predictivo}} \\ \vspace{2pt}
    {\small
        Remote / Global \ \textbullet \
        +1 (555) 019-2834 \ \textbullet \
        \href{mailto:alex.rivera@example.com}{alex.rivera@example.com} \ \textbullet \
        \href{https://alexrivera.dev}{alexrivera.dev} \ \textbullet \
        \href{https://github.com/alexrivera-dev}{github.com/alexrivera-dev} \ \textbullet \
        \href{https://linkedin.com/in/alexrivera-dev}{linkedin.com/in/alexrivera-dev}
    }
\end{center}

\vspace{-6pt}

\resumesection{Perfil Profesional}
{\small Analista de Datos y estudiante de Ciencias de la Computación en la Technical University Institute. Especializado en análisis exploratorio de datos (EDA), modelado predictivo de Machine Learning supervisado y no supervisado, ingeniería de características (Feature Engineering) y optimización de consultas SQL complejas. Experiencia práctica en preprocesamiento de datos masivos con Python (Pandas, Scikit-Learn, NumPy), diseño de pipelines ETL, segmentación con K-Means/PCA, modelos Random Forest y construcción de dashboards para toma de decisiones de negocio.}

\resumesection{Experiencia Profesional}

\textbf{Full Stack / Data-Driven Software Engineer} \hfill \textbf{Enterprise Cloud Systems} \\
\textit{Analítica de Datos de Inferencia, Optimización de Esquemas y SQL} \hfill Jun 2026 -- Actualidad
\begin{itemize}
    \item Analicé patrones de consumo y distribución estadística de payloads en llamadas a modelos de lenguaje, implementando un microservicio que redujo en un \textbf{30\% el gasto en tokens} mediante normalización de datos y compresión semántica.
    \item Diseñé y ejecuté planes de migración de datos relacionales asegurando integridad referencial, consistencia de tipos y cero pérdida de registros históricos.
    \item Optimicé consultas relacionales complejas y auditorías de datos, mitigando inyecciones SQL y mejorando tiempos de respuesta de reportes.
\end{itemize}

\vspace{1.5pt}

\textbf{Database \& Data Analyst (Contract / Freelance)} \hfill \textbf{Plataformas de Gestión \& E-Commerce} \\
\textit{Modelado Relacional, SQL Avanzado (CTEs) y Pipelines de Datos} \hfill Ene 2026 -- Actualidad
\begin{itemize}
    \item Diseñé modelos de bases de datos relacionales en PostgreSQL implementando Common Table Expressions (CTEs), procedimientos almacenados, triggers y vistas analíticas para KPIs operativos y financieros.
    \item Creé dashboards analíticos de reportería administrativa para toma de decisiones basada en datos (comportamiento de compras, inventario y retención).
    \item Implementé mecanismos de anonimización y pseudonimización de registros sensibles con HMAC para cumplimiento de privacidad de datos.
\end{itemize}

\vspace{1.5pt}

\textbf{Auxiliar de Cátedra (Teaching Assistant)} \hfill \textbf{Technical University Institute} \\
\textit{Estructuras de Datos, Algoritmos y Análisis de Complejidad} \hfill Ene 2026 -- Jun 2026
\begin{itemize}
    \item Asesoré semanalmente a 40+ estudiantes universitarios en modelado matemático, diseño de estructuras de datos y análisis de complejidad algorítmica ($O(n)$) para procesamiento eficiente de información.
\end{itemize}

\resumesection{Proyectos de Análisis de Datos y Machine Learning}

\textbf{Modelo Predictivo del Impacto de la Inteligencia Artificial en el Empleo} \textbar\ \textit{Python, Scikit-Learn, Pandas} \hfill Tech University
\begin{itemize}
    \item Lideré el análisis exploratorio (EDA) y modelado predictivo multivariado sobre métricas de riesgo laboral por automatización: tratamiento de outliers con rango intercuartílico (IQR), codificación One-Hot y escalado de variables con StandardScaler.
    \item Diseñé partición temporal de datos para evitar fuga de información (\textit{data leakage}), entrenando y comparando modelos (\textbf{Random Forest, Árboles de Decisión y K-Means}) con validación cruzada K-Fold para predecir niveles de disrupción salarial y de habilidades.
\end{itemize}

\vspace{1pt}

\textbf{Segmentación de Clientes en Retail con K-Means y PCA} \textbar\ \textit{Python, Scikit-Learn, Seaborn} \hfill \href{https://github.com/alexrivera-dev/Hoja-de-Trabajo-Kmeans}{github.com/alexrivera-dev/Hoja-de-Trabajo-Kmeans}
\begin{itemize}
    \item Construí un pipeline de segmentación de clientes no supervisado: análisis de distribuciones y correlaciones, selección de $K$ óptimo mediante método del codo y Silhouette Score, y visualización ortogonal con PCA en 2D para derivar perfiles de gasto y acciones de marketing.
\end{itemize}

\vspace{1pt}

\textbf{Brain Lab — Pipeline ETL y Motor de Búsqueda Semántica} \textbar\ \textit{Python, Vector Embeddings, BM25, SQLite}
\begin{itemize}
    \item Diseñé un pipeline ETL para ingestar y estructurar miles de interacciones conversacionales no estructuradas bajo esquemas JSON tipados.
    \item Implementé un motor de recuperación híbrido combinando búsqueda léxica BM25 y similitud vectorial, indexando la información en SQLite.
\end{itemize}

\resumesection{Habilidades Técnicas}
\begin{skilllist}
    \item \textbf{Análisis \& Manipulación de Datos:} Python (Pandas, NumPy, SciPy), Exploratory Data Analysis (EDA), Limpieza de Datos, Feature Engineering, Detección de Outliers (IQR), One-Hot Encoding, Normalización / Escalado (StandardScaler)
    \item \textbf{Machine Learning \& Modelado:} Scikit-Learn, Random Forest, Árboles de Decisión, K-Means Clustering, PCA, Regresión Lineal y Logística, Validación Cruzada (K-Fold), Métricas de Evaluación (MAE, RMSE, $R^2$, F1-Score, Curva ROC/AUC)
    \item \textbf{Bases de Datos \& SQL:} SQL Avanzado, PostgreSQL, MySQL, MariaDB, SQLite, Common Table Expressions (CTEs), Procedimientos Almacenados, Triggers, Modelado Relacional (Diagramas ER)
    \item \textbf{Visualización \& BI:} Matplotlib, Seaborn, Metabase, Figma, Excel Avanzado / Hojas de Cálculo
\end{skilllist}

\resumesection{Educación}
\textbf{Licenciatura en Ingeniería en Ciencias de la Computación} \hfill Technical University Institute \\
\textit{Cursando sexto semestre (Pensum enfocado en Sistemas, Algoritmos y Arquitectura)} \hfill 2024 -- Presente

\end{document}
